%-----------------------------------LICENSE------------------------------------%
%   This file is free software: you can redistribute it and/or                 %
%   modify it under the terms of the GNU General Public License as             %
%   published by the Free Software Foundation, either version 3 of the         %
%   License, or (at your option) any later version.                            %
%                                                                              %
%   This file is distributed in the hope that it will be useful,               %
%   but WITHOUT ANY WARRANTY; without even the implied warranty of             %
%   MERCHANTABILITY or FITNESS FOR A PARTICULAR PURPOSE.  See the              %
%   GNU General Public License for more details.                               %
%                                                                              %
%   You should have received a copy of the GNU General Public License along    %
%   with this file.  If not, see <https://www.gnu.org/licenses/>.              %
%------------------------------------------------------------------------------%
%   Author: Ryan Maguire                                                       %
%   Date:   2023/10/14                                                         %
%------------------------------------------------------------------------------%
\documentclass[a4paper,sans]{moderncv}
\usepackage[scale=0.75]{geometry}
\usepackage{multicol}
\usepackage{amssymb}

\moderncvcolor{blue}
\moderncvstyle{casual}
\name{Ryan}{Maguire}
\title{Resume}
\email{rmaguire@mit.edu}
% \phone[mobile]{+1~(555)-555-5555}

\begin{document}
    \makecvtitle
    \vspace{-8ex}
    \section{Education}
        \cventry{2012--2016}
            {B.S.}
            {University of Massachusetts - Lowell}
            {}
            {}
            {Double Major in Mathematics and Physics}
        \cventry{2016--2017}
            {M.S.}
            {University of Massachusetts - Lowell}
            {}
            {}
            {Applied Mathematics}
        \cventry{2019--2021}
            {M.A.}
            {Dartmouth College}
            {}
            {}
            {Mathematics}
        \cventry{2019--2024}
            {Ph.D.}
            {Dartmouth College}
            {}
            {}
            {Mathematics}
    \section{Experience}
        \cventry%
            {2024--Present}
                {Postdoctoral Associate}
                {MIT}
                {Cambridge, MA}
                {}
                {%
                    \begin{itemize}
                        \item
                            Instructor for the mathematics department.
                            \begin{itemize}
                                \item
                                    18.100P: Real Analysis (Spring 2025, 2026).
                            \end{itemize}
                        \item
                            Designed a vector calculus course for the
                            open learning platform.
                        \item
                            Helping to develop the VRMIT project to bring
                            VR / AR technologies into the classroom
                            (\url{https://github.com/ryanmaguire/vrmit}),
                            emphasizing vector calculus, classical mechanics,
                            and electromagnetism.
                        \item
                            Researching computational mathematics, astronomy,
                            spacetimes, and topology.
                        \item
                            Directed four PRIMES groups (2025 and 2026). The
                            subjects are:
                            \begin{itemize}
                                \item
                                    High-resolution reconstructions of
                                    Saturn's rings and particle size
                                    distributions in the F-ring.
                                \item
                                    Improving the performance of diffraction
                                    reconstruction by exploring FFT methods
                                    and polynomial interpolations for the
                                    Cassini and Voyager 2 data, and application
                                    of bifurcation methods.
                                \item
                                    Fast evaluation of the Jones polynomial.
                                \item
                                    Fast evaluation of Khovanov homology.
                            \end{itemize}
                        \item
                            Supervising research for undergraduates
                            (MIT UROP) in applied analysis,
                            mathematical physics, and astronomy.
                        \item
                            Led Directed Reading Program groups during
                            the winter sessions in 2025 and 2026.
                    \end{itemize}
                }
                \cventry%
                {2017--Present}
                {Research Assistant / Research Scientist}
                {Wellesley College}
                {Wellesley, MA}
                {}
                {%
                    Working with the Cassini Radio Science
                    Experiment and Voyager data from Saturn and Uranus
                    studying ring occultations.
                    \begin{itemize}
                        \item
                            Working in diffraction theory studying Fredholm
                            equations, Fresnel inversion, Fourier optics,
                            and signal processing.
                         \item
                            Developing open-source software in C and Python to
                            study the radio science data from the
                            Cassini mission to Saturn.
                         \item
                            Created new numerical techniques that have allowed
                            us to analyze ring structures with resolutions
                            down to 10-20 meters. Current NASA PDS data are
                            limited to 1000-meter resolutions.
                     \end{itemize}
                 }
        \cventry%
            {2021--2023}
            {Lecturer}
            {Dartmouth College}
            {Hanover, NH}
            {}
            {%
                \begin{itemize}
                    \item Math 3, Calculus (Fall 2021).
                    \item Math 54, Point-Set Topology (Summer 2022, 2023).
                    \item Math 87, Topics in Topology (Fall 2022).
                \end{itemize}
            }
        \cventry%
            {2020--2024}
            {Research Assistant}
            {Dartmouth College}
            {Hanover, NH}
            {}
            {%
                Worked in topology, knot theory, and Lorentz geometry.
                \begin{itemize}
                    \item
                        Designed and implemented algorithms for the
                        computations of knot invariants.
                    \item
                        Created large databases of knot invariants for studying
                        and experimenting with conjectures.
                    \item
                        Studied globally hyperbolic spacetimes and the
                        relationships between link invariants and causality.
                    \item
                        Wrote code to ray trace black holes in
                        various settings (Schwarzschild and de Sitter metrics).
                \end{itemize}%
             }
        \cventry%
            {2019--2021}
            {Teaching Assistant}
            {Dartmouth College}
            {Hanover, NH}
            {}
            {%
                \begin{itemize}
                    \item
                        Math 3 (Differential Calculus).
                    \item
                        Math 8 (Differential and Integral Calculus).
                    \item
                        Math 13 (Vector Calculus).
                    \item
                        Math 17 (Knot Theory and Geometry).
                    \item
                        Math 32 (The Shape of Space).
                \end{itemize}
            }
        \cventry%
            {2012--2017}
            {Research Assistant}
            {Lowell Center for Science and Space Technology}
            {Lowell, MA}
            {}
            {%
                Worked in spectroscopy, interferometry,
                LIDAR, and data analysis.
                \begin{itemize}
                    \item
                        Worked on DWEL, a dual-wavelength LIDAR that was
                        deployed in Australia, California, and Harvard Forest.
                    \item
                        Constructed and deployed HiT\&MIS. The
                        instrument was successfully deployed to Kiruna,
                        Sweden, working with Dartmouth College's BARREL
                        Mission. Performed data analysis on the airglow of
                        several events.
                    \item
                        Assisted with laboratory operations and testing for the
                        PICTURE sounding rocket.
                    \item
                        Worked on the data from the SPINR sounding rocket.
                        Developed code in IDL and worked on archiving the
                        data for MAST.
                \end{itemize}%
             }
        \cventry%
            {2015--2019}
            {Teaching Assistant}
            {University of Massachusetts - Lowell}
            {Lowell, MA}
            {}
            {%
                \begin{itemize}
                    \item
                        Management Pre-Calculus.
                    \item
                        Differential Calculus.
                    \item
                        Differential Equations.
                    \item
                        Electromagnetism (6 semesters total).
                    \item
                        Mathematics tutor for disability services.
                    \item
                        General tutoring for the physics department.
                    \item
                        General tutoring for the mathematics department.
                \end{itemize}%
            }
    \section{Professional Development}
        \begin{itemize}
            \begin{multicols}{2}
                \item ITAR Control Module, January 2013
                \item Laboratory Safety Training, January 2013
                \item DCAL Teaching Training, Winter 2019-2020
            \end{multicols}
        \end{itemize}
    \section{Computer skills}
        \makebox[2.3cm][l]{Advanced}
        \texttt{C, Python, Linux/Unix, Bash}, \LaTeX\par
        \makebox[2.3cm][l]{Intermediate}
        \texttt{C\#, C++, D, Objective-C, Fortran, Rust, Pascal,}\par
        \makebox[2.3cm][l]{\phantom{Intermediate}}
        \texttt{Go, IDL, Swift, Java, JavaScript, SageMath}\par
        \makebox[2.3cm][l]{Elementary}
        \texttt{Mathematica, MATLAB, Ruby, Julia, Kotlin, R, TypeScript}\par
\end{document}
